\chapter{Introduction}
A Compact Heat Exchanger is a class of heat exchanger designed to maximize heat transfer efficiency within a minimal volume, characterized by a high surface area-to-volume ratio, typically exceeding 700 m²/m³. It facilitates efficient thermal energy exchange between two fluids (gas or liquid) with closely spaced, often intricately structured, heat transfer surfaces, such as fins, tubes, or plates, to enhance heat transfer while minimizing size and weight. These exchangers are commonly used in aerospace, automotive, HVAC applications and industrial processes where space and efficiency are critical.

The seminal text on the subject by William Kays and Lou London, entitled "Compact Heat Exchangers," emphasizes some senior considerations when designing heat exchangers in general. For example, while the purpose of a heat exchanger is to transfer heat from a hot fluid to a colder fluid and raise its temperature where fluid is either a gas or liquid. In order to accomplish this transfer of energy, the cold and hot fluids require mechanical pumping power to overcome the friction from moving a fluid through heat transfer passages.

\section{Estimating Power Losses}

